% 阶段报告四导读。编译：xelatex summary_plain_zh.tex（两遍）。
\documentclass[12pt,a4paper]{article}
\linespread{1.10}
\usepackage[margin=0.9in]{geometry}
\usepackage{amsmath,booktabs,tabularx,graphicx,xeCJK}
\setCJKmainfont{FandolSong-Regular}[Extension=.otf,BoldFont=FandolSong-Bold]
\setCJKsansfont{FandolHei-Regular}[Extension=.otf,BoldFont=FandolHei-Bold]
\usepackage[hidelinks]{hyperref}
\renewcommand{\figurename}{图}
\renewcommand{\tablename}{表}
\emergencystretch=3em
\setlength{\parskip}{4pt}
\title{\bfseries 尺度增益的来源、Dead 专家分支的关闭与推理成本\\[7pt]
\large 阶段报告四导读}
\author{课程项目 bdd26}
\date{2026 年 10 月 8 日}
\begin{document}
\maketitle

\section*{核心结论}

本阶段不以求更高的测试分数为目标，而是把上一份报告留下的三个问题逐一推进到明确的答案：

\begin{itemize}
\item \textbf{两倍尺度的 Dead 增益，主要不是训练门槛带来的。}模型训练沿用了官方实现的一个 30 像素面积门槛：面积小于该值的核，其辅助 HV 几何监督取值全为零。两倍尺度下，同一门槛实际对应更小的原图面积，相当于一个未受控的变量。对照实验显示，在原尺度上单独放宽这个门槛，只能复现两倍尺度 Dead 增益的 \textbf{9.59\%}。
\item \textbf{Dead 专属专家分支这一研究线，在投入训练算力之前即按事先规则关闭。}四道预注册闸门里，三道通过、一道失败：训练好的模型上，Dead 像素与常见类像素的梯度\textbf{方向高度一致}，并不冲突——“靠专家分支把 Dead 从共同训练里解耦出来”这一前提不成立。
\item \textbf{两套模型的完整推理成本实测约为单模型的 6.09 倍。}按事先的规则，EF-P2 从“候选部署方案”调整为“机制证据与教师候选”，后续任何部署层面的表述都须同时给出这一成本。
\end{itemize}

\section*{训练目标中的面积门槛：一个与尺度耦合的变量}

上一阶段发现：把图像放大到 $512\times512$ 再配套训练，能比原尺度多检出一些原本漏检的 Dead 核，但整体分割质量略有下降。放大同时改变了许多因素，收益不能笼统归因于“分辨率”。

复盘代码时发现一处具体差异：训练目标里的 HV 通道（为每个核编码“到中心的距离与方向”的辅助地图，解码时用于把相邻的核分开）沿用了官方的 30 工作像素门槛——面积小于该值的核，这张地图取值全为零。同一个 30 像素门槛，在两倍尺度上只相当于原图 7.5 个像素。于是一个 16 像素的小核：原尺度下地图全零，两倍尺度下有完整的地图。\textbf{两种尺度给小核的几何监督，在训练开始前就不一样}。

该门槛是官方实现的既有行为，并非程序缺陷；修改它属于实验干预，需要以对照实验评估。

\section*{四个对照臂：门槛解释的增益不足一成}

实验设计成四个臂：原尺度/两倍尺度，各配“默认门槛 30”和“互换门槛”（原尺度降到 8、两倍尺度升到 120），后者让两边的“原生面积支持”基本配平。每个臂两个种子，共 8 次训练，全部在验证折上评估——\textbf{这一阶段全程未使用测试折}。

三个关键读数：

\begin{itemize}
\item 原尺度把门槛从 30 降到 8：Dead 增益在两个种子间\textbf{一正一负}，整体指标还略降。降低门槛未能改善小核的检出，也不能替代第二套模型。
\item 两倍尺度内部把门槛从 120 降到 30：官方口径的 Dead 指标在两个种子上均上升（平均 $+0.0098$），但置信区间跨零，属于“候选信号”，不能当作稳健结论。
\item 配平支持后直接比尺度：Dead 增益和 bPQ 代价\textbf{都还在}——而且更严格的 strict Dead 口径在全部四个种子对上都是负的。尺度的主体效应在别处（采样、特征、训练日程都还没排除）。
\end{itemize}

对象级统计进一步支持这一判断：门槛在两倍尺度内几乎不改变“内部完整 Dead 核被匹配上的数量”（526.0 对 526.5，两个种子总共多 1 个）；也就是说，$+0.0098$ 并非来自多检出的核，更可能来自掩码形状和误报结构的变化。

\section*{Dead 专家分支：四道预注册闸门检验与路线关闭}

另一条路线更具探索性：在模型里加一个 Dead 专属分支（自带检测几何头，只在含 Dead 的图上计算损失），希望它在不降低整体指标的前提下找回漏检的 Dead。团队没有直接开始训练，而是先在预注册规格里写好\textbf{四道训练前闸门}——全部通过才启动训练，任一失败即关闭整条研究线：

\begin{center}
\begin{tabular}{lll}\toprule
闸门 & 检验什么 & 结果\\\midrule
A & 门槛能解释多少尺度增益 & 通过（解释 9.59\%，未达 50\% 的换基线条件）\\
C & 仅调类别阈值的效果上限 & 通过（峰值 $+0.002$，未达 $+0.005$ 的关闭线）\\
D & 漏检 Dead 的收益上限 & 通过（oracle 上界 $+0.156$）\\
B & Dead 与常见类梯度是否冲突 & \textbf{失败（详见下文）}\\\bottomrule
\end{tabular}
\end{center}

闸门 B 的过程本身值得记录。第一次运行按预冻结计划“顺序取前 256 张图”，结果这 256 张全部来自乳腺组织、其中仅 1 张含 Dead——64 个批次里 63 个的 Dead 损失恒为零，未测到任何梯度几何信息。这不是“失败”，而是\textbf{无效测量（VOID）}：问题出在取样设计，不在探针。修正方案（改为对 65 张 Dead 阳性图分层取样、分母仅计有效批次）先登记进台账、再补齐测试，经确认后才重跑。

重跑结果明确：64 个 Dead 阳性批次上，解码器共享参数的复合梯度余弦\textbf{平均 $+0.749$、最小值 $+0.594$、64 批中 0 批为负}。余弦为负才被判定为“冲突”；全部落在正半轴，说明优化 Dead 像素与优化常见类像素把参数往\textbf{同一个方向}推。既然不冲突，“用专家分支把 Dead 解耦出来”所要解决的前提问题就不存在了。

按事先写进规格的止损规则，这条线当天关闭：冒烟轮、开发轮、确认轮全部取消。综合来看，这条线留下一组相互印证的机制证据：Dead 短板是\textbf{检测召回}问题——重加权与平衡采样均已尝试且无效、调阈值的收益上限为 $+0.002$、专家解耦的前提不成立、而将漏检的 Dead 全部检出（oracle 上界）对应 $+0.156$。这些结果指向“新的检测监督或数据”，而不是在现有监督上重新分配容量。

\section*{推理成本实测}

最后是推理成本实测。固定协议下：单模型 x1 每图 15.5 毫秒；两倍尺度模型 52.9 毫秒；两套预测的 CPU 融合每图另需 26 毫秒——\textbf{合计 94.4 毫秒/图，是单模型的 6.09 倍}。作为参照，单模型自己做测试时增强（TTA）的成本也达到无 TTA 推理的 4.67 倍，但 TTA 并不带来新的 Dead 检出。

绝对毫秒数在共享 GPU 上随负载波动（系统平均负载 loadavg 在 19--46 之间时，x1 每图耗时在 10.7--15.5 毫秒间变化），但\textbf{臂间比值}在所有测量窗口里都稳定在 6.1--6.3。按预注册规则：不具备性能--成本优势，EF-P2 就定位为机制证据与教师候选，而非部署方案。

\section*{一致性审计：文档、代码与实验产物互核}

收尾阶段进行了一次全仓库一致性审计：文档、代码、实验产物三方互核。结论是主要结果全部复现，同时发现并更正了十几处转录级小错——例如 CellViT-UNI 的三划分 bPQ 均值，旧记录的 0.6654 是“先舍入再平均”造成的，按原始汇总精确平均应为 \textbf{0.6653}。两个在两台机器上都丢失的产物从留存的确定性预测重建，数值与历史记录逐位一致，并用哈希侧车文件绑定来源。工程上还增加了“测试折守卫”：开发类脚本此后会被阻止读取所属划分的测试折。审计后全量测试 271 项通过。

\section*{下一步工作：标注与数据侧的不确定性}

机制与成本的问题已经明确，剩下的关键不确定性在\textbf{标注与数据}一侧：EF-P2 补进来的候选里，那些“没匹配上”的属于碎片、形状偏差，还是标注本身的遗漏？下一轮首选 E1 盲审——四组候选（x1 接受的、x2 独有且被接受的、x2 独有被拒的、随机背景）交由独立评审逐个判读，允许回答“当前视野无法判定”。如果盲审支持“有一部分是真实的漏标核”，再采用训练内交叉拟合的效用学习；如果主要是碎片，重心就转向可见性与边界机制。两条路线的工具（分层蒙太奇、梯度探针、oracle 上界）这一阶段都已备好。

\vspace{1em}
\noindent\small 完整数字、置信区间、闸门台账与复现说明见同目录下的
\texttt{stage\_report\_4\_zh.pdf}；英文版为 \texttt{stage\_report\_4.pdf}。
\end{document}
